\subsection{Непреревность сложной функции, теорема о сохранении знака.}

\subsubsection*{Теорема (О сохранении знака).}
Если $f(x)$ непрерывна в $a$ и $f(a) \neq 0$, то существует окрестность $U_\delta(a)$, в которой $\operatorname{sgn} f(x) = \operatorname{sgn} f(a)$.

\begin{tcolorbox}[title=Доказательство]
	Пусть $f(a) > 0$. Возьмем $\varepsilon = \frac{f(a)}{2} > 0$.
	По определению непрерывности $\exists \delta > 0$, такое что $\forall x \in U_\delta(a)$:
	\[ |f(x) - f(a)| < \frac{f(a)}{2} \Leftrightarrow -\frac{f(a)}{2} < f(x) - f(a) < \frac{f(a)}{2} \]
	\[ \frac{f(a)}{2} < f(x) < \frac{3f(a)}{2} \]
	Так как $\frac{f(a)}{2} > 0$, то $f(x) > 0$ во всей окрестности. (Для $f(a) < 0$ аналогично).
\end{tcolorbox}

\subsubsection*{Непрерывность сложной функции (суперпозиции)}

\textbf{Теорема.}
Пусть:
\begin{enumerate}
	\item $x = \phi(t)$ непрерывна в точке $\tau$ и $\phi(\tau) = a$.
	\item $y = f(x)$ непрерывна в точке $a$.
\end{enumerate}
Тогда сложная функция $F(t) = f(\phi(t))$ непрерывна в точке $\tau$.

\begin{tcolorbox}[title=Доказательство]
	Запишем определения непрерывности для обеих функций:
	\begin{enumerate}
		\item Для $f(x)$ в точке $a$:
		\[ \forall \varepsilon > 0 \; \exists \sigma > 0: \forall x \in U_a \; (|x - a| < \sigma \Rightarrow |f(x) - f(a)| < \varepsilon) \]
		
		\item Для $\phi(t)$ в точке $\tau$ (в качестве "эпсилон" берем найденное $\sigma$):
		\[ \text{Для } \sigma > 0 \; \exists \delta > 0: \forall t \in U_\tau \; (|t - \tau| < \delta \Rightarrow |\phi(t) - \phi(\tau)| < \sigma) \]
	\end{enumerate}
	
	\textbf{Объединяем:}
	Возьмем произвольное $\varepsilon > 0$. Найдем $\sigma(\varepsilon)$ из п.1. Затем найдем $\delta(\sigma)$ из п.2.
	Тогда:
	\[ |t - \tau| < \delta \implies \underbrace{|\phi(t) - \phi(\tau)|}_{=|x - a|} < \sigma \implies |f(\phi(t)) - f(\phi(\tau))| < \varepsilon \]
	Что и означает непрерывность композиции.
\end{tcolorbox}

\begin{tcolorbox}[title=Пояснение про области определения]
	Чтобы сложная функция $f(\phi(t))$ существовала, значения внутренней функции $x = \phi(t)$ должны попадать в область определения внешней функции $f(x)$.
	
	В теореме предполагается, что $f(x)$ определена в некоторой окрестности $U_a$.
	Так как $\phi(t)$ непрерывна и $\phi(\tau) = a$, то значения $\phi(t)$ при $t$, близких к $\tau$, будут как угодно близки к $a$.
	
	Значит, найдется такая маленькая окрестность точки $\tau$, что все значения $\phi(t)$ из неё попадут внутрь области определения $f(x)$ (в окрестность $U_a$).
	Без этого условия запись $f(\phi(t))$ просто не имела бы смысла.
\end{tcolorbox}



